\section{引言}

%纳米硒作为一种新型纳米材料，在中西医结合领域具有广阔的应用前景 \cite{文献1}。
%然而，传统的体内外毒性评价实验耗时耗力，且难以覆盖复杂的理化参数组合。

%近年来，深度学习在计算毒理学中展现出巨大潜力。
%目前的研究主要集中在有机小分子，针对纳米材料（尤其是纳米硒）的预测模型尚不成熟。

%本研究旨在建立一个多参数输入的深度学习模型，通过分析纳米硒的粒径、浓度、表面修饰等特征，实现对其生物毒性的快速、准确预测。

纳米硒是一类具有优异生物相容性和多重生物活性的纳米材料，在中西医结合领域具有广阔的应用前景 \cite{liu2012senps,wang2022tcmp}。
然而，纳米硒的生物效应高度依赖于粒径、浓度、表面修饰等理化参数，并在某些情况下产生细胞毒性，从而显著影响其生物安全性。


近年来，深度学习在高维特征挖掘和非线性关系建模方面展现出显著优势，为计算毒理学提供了新的研究范式 \cite{guo2023toxicity}。
然而，现有的预测模型主要面向有机小分子，对于纳米材料尤其是纳米硒的毒性预测研究仍较有限，相关建模方法尚不成熟。

针对这一问题，本研究构建了一种多特征融合的深度学习模型，以纳米硒的粒径、浓度和表面修饰等关键理化参数为输入，旨在实现其生物毒性的快速、准确预测，为纳米硒的生物安全性评估和合理应用提供新的工具与思路。